\newpage

\section{Trabajo a futuro}

Para la siguiente fase del proyecto, se tiene planificada la ejecución de los Sprints 3 y 4 de la metodología Scrum establecida \cite{sutherland2020scrum}. Estas etapas contemplan el desarrollo técnico, la integración de componentes y la validación final del prototipo en un entorno real. Las actividades se desglosan en las siguientes áreas de acción:

    \textbf{Optimización del modelo y mitigación de riesgos}: Una prioridad inmediata es el incremento del volumen del conjunto de datos de otras temporadas del año. Se proyecta realizar campañas intensivas de captura en campo para alcanzar un umbral de entre 10,000 y 20,000 imágenes por clase, lo que permitirá reducir la brecha de generalización observada en las curvas de entrenamiento de las pruebas preliminares \cite{goodfellow2016deep}. Estas capturas se realizarán bajo diversas condiciones de iluminación natural y ángulos para mitigar el riesgo de sesgo de dominio. Asimismo, se implementará un proceso de \textit{fine-tuning} sobre la arquitectura EfficientNet-B3 \cite{Tan2019EfficientNetRM}, priorizando la mejora de la métrica de sensibilidad (\textit{recall}), la cual es fundamental para minimizar los falsos negativos en la detección temprana de estas plagas \cite{forestal2025insectos}.

    \textbf{Implementación de la arquitectura Cloud y Offline-First}: Se llevará a cabo el despliegue de la infraestructura \textit{serverless} en Amazon Web Services (AWS) \cite{cloudplexo_serverless, northflank_comparison}. El núcleo del sistema será una API REST desarrollada en Micronaut y Java 21, la cual gestionará la lógica de negocio mediante el enfoque de \textit{Clean Architecture} y administrará la persistencia políglota en Amazon RDS (PostgreSQL) \cite{aws_rds_pricing} y Amazon S3. En el cliente móvil, se consolidará la estrategia \textit{Offline-First}, asegurando que el modelo de clasificación \cite{omahony2019deep_learning_vs_traditional_cv} y el procesamiento de color con OpenCV operen de forma concurrente en el dispositivo. La sincronización de datos se programará mediante procesos en segundo plano que detecten de forma automática la disponibilidad de red para subir los registros locales alojados en SQLite hacia la nube.

    \textbf{Escalabilidad y proyección a ambiente productivo}: El sistema ha sido diseñado desde su arquitectura base para escalar hacia un entorno de producción. A futuro, se contempla la transición desde los límites de la capa gratuita de AWS \cite{aws_free_tier} hacia instancias de mayor capacidad, como \textit{db.t3.medium} para base de datos \cite{aws_calculator_prototipo, aws_calculator}, y la implementación de AWS WAF para la protección contra amenazas externas en la red. Debido a la arquitectura modular empleada, el sistema podrá adaptarse a los cambios climaticos dada la naturaleza del dataset, que se compone de imagenes tomadas en todas las temporadas del año \cite{delval_saenzromero2017descortezadores, vega2022mongabay} mediante el reentrenamiento del modelo en Amazon SageMaker, permitiendo que esta solución sea replicable en otros bosques y pulmones urbanos del país \cite{forestal2025bosques}.

    \textbf{Validación integral y documentación}: Finalmente, se ejecutará el plan general de pruebas funcionales, de integración y de sistema diseñado en la primera etapa. Se planea involucrar a personal forestal operativo en sesiones de prueba, con el objetivo de identificar áreas de mejora en la usabilidad tanto de la aplicación Android como de la plataforma web.
